\documentclass[11pt,a4paper]{article}
\usepackage[margin=27mm]{geometry}
\usepackage[T1]{fontenc}
\usepackage{lmodern,microtype,amsmath,amssymb,amsthm,booktabs,array}
\usepackage[hidelinks]{hyperref}
\hypersetup{pdftitle={An Ant on the Clock: Exponential Escape Bounds for Langton's Ant},pdfauthor={Hao Ke}}
\usepackage{fancyhdr}
\pagestyle{fancy}\fancyhf{}\fancyhead[L]{\small An Ant on the Clock}\fancyfoot[C]{\thepage}
\setlength{\headheight}{14pt}
\newtheorem{theorem}{Theorem}[section]
\newtheorem{lemma}[theorem]{Lemma}
\newtheorem{corollary}[theorem]{Corollary}
\theoremstyle{remark}\newtheorem{remark}[theorem]{Remark}
\newcommand{\Z}{\mathbb Z}\newcommand{\N}{\mathbb N}\newcommand{\Q}{\mathbb Q}
\newcommand{\eps}{\varepsilon}\newcommand{\symdiff}{\mathbin{\triangle}}
\title{An Ant on the Clock:\\Exponential Escape Bounds for Langton's Ant}
\author{Hao Ke\\Independent researcher\\\texttt{i@kehao.me}}
\date{13 October 2026}
\begin{document}
\maketitle
\begin{abstract}
We prove that the classical two-colour Langton ant exits every square of side
$n$ within $8n3^n+1$ updates, uniformly over all initial internal colourings,
positions and headings. This improves the asymptotic order of Etse's factorial
square bound to exponential in the side length. The proof constructs a family
of rational clocks that increase strictly at every update: a bound on the
clock's range inside the square then gives a bound on the time spent there.
On complete finite-support states, a fixed clock increases by at least $2/5$
per update and yields decidable reachability, two-sided orbit equivalence and
a computable orbit normal form. Exact finite checks and Lean proofs of the
local algebra and telescoping lemmas accompany the analytic argument.
\end{abstract}

\section{Model and results}
A state is $X=(B,p,h)$, where $B\subset\Z^2$ is finite, $p\in\Z^2$, and
$h\in\{N,E,S,W\}$. Headings are encoded cyclically by $0,1,2,3$.
Let $e_h$ be the corresponding unit vector. A state at time $t$ is observed
\emph{before} update $t$. The literal update is
\begin{equation}\label{eq:rule}
 h'=h+1-2\mathbf1_B(p)\pmod4,\qquad
 B'=B\symdiff\{p\},\qquad p'=p+e_{h'}.
\end{equation}
Thus white gives a right turn, black a left turn, and the read cell flips
before the ant moves.

There are two distinct questions in this note. The first asks how long the ant
can remain in a specified finite square before its first exit. The second asks,
for two \emph{specified complete states} $X,Y$, whether $F^tX=Y$ for some
integer $t\ge0$. ``Complete'' includes every black cell, not just a neighbourhood
of the ant or a retained active motif.

Etse~\cite{Etse2026} studies first escape from finite connected domains and proves,
in particular, a square bound $(n+1)!$, alongside bounds for fixed-height
rectangles. Our main quantitative result is the uniform bound
$T_{K_n}\le8n3^n+1$ for a square $K_n$ of $n$ cells per side
(Corollary~\ref{cor:exp}). The estimate includes the exit move and holds for
every internal colouring and pose, independently of exterior colours.
The same construction gives algorithms whose inputs are two complete finite
black sets and ant poses, complementing the circuit-prediction and
infinite-background questions studied in~\cite{Gajardo2002}.

The common mechanism is a rational clock with a positive minimum increment.
Inside a square, its finite range limits the number of updates before exit;
for a specified complete target state, the difference of clock values bounds
the time at which that target can be reached. We first construct the clock
and prove its increment identity, then estimate its range on a square and
choose parameters to obtain the exponential bound. Finally, we fix one clock
to derive complete-state reachability and a normal form for two-sided orbits.

\section{An explicit strict rational clock}
Define
\[
 \chi(x,y)=(-1)^{x+y},\qquad
 \sigma(X)=\begin{cases}\chi(p),&h=N,S,\\-\chi(p),&h=E,W,\end{cases}
\]
and put $\eps(h)=+1$ for $h=N,E$, and $\eps(h)=-1$ for $h=S,W$.
Each update changes both the lattice parity and the heading axis, so $\sigma$
is invariant under $F$.

The construction uses the diagonal coordinate $x+y$, which changes by
exactly $+1$ or $-1$ on every move. We assign symmetric exponential weights
to the two directions along this coordinate. A colour flip changes one term
of the weighted field sum; a heading-dependent pose term compensates for that
change so that their combined increment is positive in every turn case.
The invariant $\sigma$ aligns the field signs with the ant's heading axis.

Fix $c\in\Z$ and a rational $r>3$. The parameter $c$ centres the diagonal
coordinate, while $r$ controls the growth of the weights. Write
$s(q)=q_x+q_y-c$ and set
\begin{align}
 f_r(s)&=r^s+r^{-s},\\
 g_r(s,+1)&=\frac{2}{r+1}(r^s-r^{1-s}),\\
 g_r(s,-1)&=\frac{2}{r+1}(r^{-s}-r^{s+1}),\\
 \Phi_{r,c}(X)&=-\sigma(X)\sum_{q\in B}\chi(q)f_r(s(q))
                    +g_r(s(p),\eps(h)).\label{eq:clock}
\end{align}
Finite support makes \eqref{eq:clock} a computable rational number, including
when coordinates are negative.

\begin{lemma}[Increment identity]\label{lem:increment}
Let $U=r^{s(p)}$, $V=r^{-s(p)}$, and
\[
 d_r=\frac{r-3}{r+1},\qquad E_r=\frac{5r+1}{r+1}.
\]
If $e=\eps(h)$ and $e'=\eps(h')$, then
\[
\begin{array}{c|c}
 (e,e')&\Phi_{r,c}(FX)-\Phi_{r,c}(X)\\\hline
 (+,+)\text{ or }(-,-)&d_r(U+V)\\
 (+,-)&d_rU+E_rV\\
 (-,+)&E_rU+d_rV
\end{array}
\]
In particular $\Phi_{r,c}(FX)-\Phi_{r,c}(X)\ge2d_r>0$.
\end{lemma}
\begin{proof}
Let $\delta=+1$ on a white read and $\delta=-1$ on a black read. From the
four headings and two colours in \eqref{eq:rule},
\[
 \delta\sigma(X)\chi(p)=ee',\qquad s(p')=s(p)+e'.
\]
The black-field contribution to the increment is therefore
$-ee'(U+V)$. Subtracting the old pose term from the new one gives
\[
 -ee'(U+V)+g_r(s(p)+e',e')-g_r(s(p),e).
\]
Substitution of the two definitions of $g_r$ gives the displayed table.
For example, for $(e,e')=(+,+)$, the pose increment is
$2(r-1)(U+V)/(r+1)$, which yields $d_r(U+V)$ after the field increment.
For completeness, after multiplying by $r+1$, the other three cases are
\begin{align*}
(-,-):\quad&-(r+1)(U+V)+2(rV-U)-2(V-rU),\\
(+,-):\quad& (r+1)(U+V)+2(rV-U)-2(U-rV),\\
(-,+):\quad& (r+1)(U+V)+2(rU-V)-2(V-rU).
\end{align*}
Expansion gives the corresponding polynomial identities. Since $UV=1$, $U+V\ge2$;
also $E_r\ge d_r>0$. These give the lower bound in every case.
\end{proof}

The parameter family will let us balance the clock's increment against its
range on a square. For complete-state orbit algorithms, one fixed member
suffices: choosing $r=4$ and $c=0$ gives the following uniform time estimate.

\begin{theorem}[Strict clock]\label{thm:clock}
The function $\Phi=\Phi_{4,0}$ is computable on every complete finite-support
state and satisfies
\begin{equation}\label{eq:time}
 \Phi(F^tX)-\Phi(X)\ge\frac25t\qquad(t\in\N).
\end{equation}
\end{theorem}
\begin{proof}
Apply Lemma~\ref{lem:increment} with $r=4$ and sum the increments.
\end{proof}

\section{Exponential first exit from a square}
To estimate confinement time, we use the freedom to centre the weights on
the square and choose their base as a function of its side length. We first
retain $r>3$ in the range estimate, then specialize to $r=3+3/n$.

Let $K_n=\{0,\ldots,n-1\}^2$, with $n\ge1$. Start with $p_0\in K_n$ and
any heading. Put
\[
 T_{K_n}=\min\{t\ge1:p_t\notin K_n\},
\]
when the set is nonempty. The exit move itself is counted. Coordinates of a
translated square are obtained simply by translating the whole argument.

For this section, exterior colours may even be arbitrary: define the potential
by summing over $B\cap K_n$ only. Every read up to the exit is inside $K_n$,
so Lemma~\ref{lem:increment} remains valid for internal-to-internal updates.
Thus the first-exit estimate depends only on the square's internal colouring.

\begin{theorem}[Parameterized square bound]\label{thm:square}
For every rational $r>3$,
\begin{equation}\label{eq:square}
 T_{K_n}\le1+\left\lfloor\frac{r+1}{r-3}
 \left[\frac{r^{1-n}(r^n-1)^2}{(r-1)^2}
             +r^{n-1}-r^{1-n}\right]\right\rfloor.
\end{equation}
In particular $T_{K_n}$ always exists.
\end{theorem}
\begin{proof}
Use $c=n-1$ and let $R=n-1$. For all positions in $K_n$ one has $-R\le s\le R$.
With fixed $\sigma$, the range of the black-field sum over all internal
colourings is at most
\begin{align}
 F_n&=\sum_{q\in K_n}(r^{s(q)}+r^{-s(q)})\\
 &=2r^{-R}\left(\frac{r^n-1}{r-1}\right)^2.\label{eq:fieldrange}
\end{align}
Each coefficient may be either present or absent; the difference between its
largest and smallest possible contribution is its absolute value. Summing those
ranges proves the first line. Reflection of the square exchanges the two
exponential sums and proves the second.

The function $g_r(s,+1)$ increases with $s$, and
$g_r(s,-1)=g_r(-s,+1)$. Consequently the range of the pose correction is at most
\begin{align}
 G_n&=g_r(R,+1)-g_r(-R,+1)\\
 &=2(r^R-r^{-R}).\label{eq:poserange}
\end{align}
Thus all internal values of the potential lie in an interval of length
$F_n+G_n$. If the ant never exited, its potential would increase by at least
$2d_r$ at every update while staying in this bounded interval, a contradiction.
Hence the first exit is finite.

States at times $0,\ldots,T_{K_n}-1$ are internal. The $T_{K_n}-1$ updates
between these states give
\[
 2d_r(T_{K_n}-1)\le F_n+G_n.
\]
Divide, use \eqref{eq:fieldrange}--\eqref{eq:poserange}, and take the integer
part to obtain \eqref{eq:square}. The additional unit counts the exit update.
\end{proof}

\begin{corollary}[Exponential bound in side length]\label{cor:exp}
For every initial colouring and internal pose,
\[
 \boxed{T_{K_n}\le8n3^n+1.}
\]
\end{corollary}
\begin{proof}
Choose $r=3+3/n$. The factor $(r+1)/(r-3)=(4n+3)/3$ is at most $7n/3$.
Since $r>3$,
\[
 \frac{r^{1-n}(r^n-1)^2}{(r-1)^2}
 \le\frac{r^{n+1}}{(r-1)^2}\le\frac94r^{n-1}.
\]
The other two terms in brackets in \eqref{eq:square} have sum at most $r^{n-1}$.
The binomial theorem gives
$(1+1/n)^n\le\sum_{k=0}^n1/k!<3$: the terms for $k=0,1$ sum to two,
and for $k\ge2$, $k!\ge2^{k-1}$ bounds the remaining finite sum by a
strictly subunit geometric sum. It follows that $r^{n-1}<3^n$. The expression inside the floor
in \eqref{eq:square} is less than
\[
 \frac{7n}{3}\,\frac{13}{4}\,3^n
 =\frac{91}{12}n3^n<8n3^n.
\]
The stated convenient integer bound follows.
\end{proof}

\begin{remark}[Comparison of square bounds]
Etse's Theorem~5~\cite{Etse2026} gives $T_{K_n}\le(n+1)!$. Corollary~\ref{cor:exp}
gives the asymptotically smaller order $O(n3^n)$ in the side length $n$.
For an individual square, the minimum of Etse's bound and
\eqref{eq:square} gives a useful combined estimate. The quantity bounded
throughout is the first exit time, including the move across the boundary.
\end{remark}

\begin{center}
\begin{tabular}{r r r}
\toprule
Side $n$ & Bound \eqref{eq:square}, $r=3+3/n$ & $(n+1)!$\\
\midrule
3 & 218 & 24\\
6 & 13,982 & 5,040\\
8 & 175,691 & 362,880\\
12 & 22,419,218 & 6,227,020,800\\
16 & 2,484,936,650 & 355,687,428,096,000\\
\bottomrule
\end{tabular}
\end{center}
The table compares upper bounds at selected side lengths: the factorial bound
is smaller in the first two rows, and the clock bound is smaller in the last three.

\section{Complete-state orbit algorithms}
Fix the single clock $\Phi=\Phi_{4,0}$ throughout this section.
Here the comparison is between complete states, so the same fixed clock is
used for every input. Its uniform increment converts a target clock value
into a finite simulation budget.

\begin{theorem}[Forward reachability]\label{thm:reach}
Given complete finite-support states $X,Y$, it is decidable whether
$F^tX=Y$ for some $t\ge0$. If such a time exists it is unique and satisfies
\[
 t\le N(X,Y):=\left\lfloor\frac52(\Phi(Y)-\Phi(X))\right\rfloor.
\]
\end{theorem}
\begin{proof}
Different $\sigma$ values or a negative clock difference rule out a match.
Otherwise compute the displayed integer and simulate through times $0$ to
$N(X,Y)$, inclusive, comparing the entire black set, position and heading.
Inequality~\eqref{eq:time} rules out every later match. Strict increase also
rules out two different matching times. Each state at finite time is finite,
and all rational arithmetic and full-state comparisons are effective.
\end{proof}

The inverse update is equally explicit. Given $(B,p,h)$, put $q=p-e_h$.
If $q\in B$, the preceding update was white-to-black, so the old heading was
$h-1$. If $q\notin B$, it was black-to-white, so the old heading was $h+1$.
In either case flip $q$ and put the predecessor at $q$. This defines a computable
bijection on the set $\mathcal X$ of complete finite-support states.

\begin{corollary}[Orbit equivalence]
The relation $\exists t\in\Z\;F^tX=Y$ is decidable, with a computable signed
time offset whenever it holds.
\end{corollary}
\begin{proof}
If the states differ, their time order on an orbit must agree with the strict
order of their clock values. Apply Theorem~\ref{thm:reach} in that order;
equal clock values allow a match only for identical states.
\end{proof}

\begin{theorem}[Computable orbit section]\label{thm:section}
Define the decidable set
\[
 \mathcal S=\{C\in\mathcal X:\Phi(C)\ge0>\Phi(F^{-1}C)\}.
\]
Every two-sided orbit meets $\mathcal S$ exactly once. Its representative $C(X)$
and the unique integer $k(X)$ with $X=F^{k(X)}C(X)$ are computable. The map
$X\mapsto(C(X),k(X))$ is a computable bijection between $\mathcal X$ and
$\mathcal S\times\Z$ under which $F$ acts by $(C,k)\mapsto(C,k+1)$.
\end{theorem}
\begin{proof}
Forward values tend to $+\infty$ by \eqref{eq:time}. Applying the same inequality
to $F^{-t}X$ shows that backward values tend to $-\infty$. Strict increase
therefore yields exactly one crossing from negative to nonnegative values.
Search forward when $\Phi(X)<0$ and backward otherwise, counting the signed
number of moves. A coarse bound for the moves and final predecessor check is
$\lceil(5/2)|\Phi(X)|\rceil+1$. The inverse of the displayed correspondence
computes $F^kC$ using forward or inverse updates as appropriate.
\end{proof}

The normal form separates the position along an orbit from its orbit label
$C\in\mathcal S$. A row-parity invariant shows that infinitely many distinct
labels are needed.

\begin{lemma}[Infinitely many orbit representatives]
The section $\mathcal S$ is infinite.
\end{lemma}
\begin{proof}
For each integer row $j$, the bit
\[
 I_j(X)=|B\cap(\Z\times\{j\})|
        +\mathbf1_{\{p_y=j,\ h\in\{E,W\}\}}\pmod2
\]
is invariant. On a vertical-to-horizontal turn the new head stays in the
read row and its indicator cancels that row's colour flip. On a
horizontal-to-vertical turn the old head indicator cancels the flip.
All other rows are unchanged. States with the head at the origin facing north
and sole black cell $(0,j)$ have distinct invariant vectors $(I_i)_{i\in\Z}$.
They therefore lie on distinct orbits, each with its own section representative.
\end{proof}

\paragraph{Computational cost.}
The clock gives explicit finite simulation budgets for the orbit algorithms.
With sparse binary coordinates, exact arithmetic may require exponential
space in the coordinate bit length: the expanded rational representation of
$4^x$ has length proportional to $|x|$. The algorithms establish decidability
with this representation-dependent cost.

\section{Verification and reproducibility}
The main theorem combines exact local identities with a finite geometric-series
range estimate. The supplementary artifact supplies two independently written
steppers and exact rational arithmetic. The retained baseline replay checks
4,704 local cases, 18,696 states in squares through side three, and 100 bound
parameters. A new integer-vector implementation checks 9,720 local cases,
including negative coordinates, remote black cells, and bases arbitrarily close
to three. A separate bit-mask functional graph exhausts all 4,213,000
colour/pose states of squares of sides one through four. Their maximum exit
times are $1,6,17,46$. Every graph transition is evaluated once, cycles are
explicitly rejected, and a maximizing witness for each size is replayed with
the integer-vector stepper. Direct field sums are compared with the geometric
formula in 60 cases; the coarse bound is checked at 1,000 side lengths.
These checks exercise the local identities and small-square dynamics alongside
the proofs of Lemma~\ref{lem:increment} and Theorem~\ref{thm:square}.

The artifact also includes kernel-checked Lean 4.30.0 proofs of all four
polynomial increment identities after clearing the positive denominator,
$U+V\ge2$ when $U>0$ and $UV=1$, the scaled increment lower bound, and the
finite telescoping/interval argument. Their printed dependencies are only
\texttt{propext}, \texttt{Classical.choice}, and \texttt{Quot.sound}.
The formalized components are the local algebra and telescoping lemmas.
The connection to the literal ant, integer-power identities, finite-field
range, and $8n3^n+1$ specialization are established by the analytic proofs above.

\section{Conclusion}
A strict rational clock connects local ant updates to quantitative control of
complete trajectories. On a square, exponential weights give a finite range
whose comparison with the minimum increment yields $T_{K_n}\le8n3^n+1$.
On complete finite-support states, the same increment bounds the time to a
specified target, while reversibility gives a unique zero-crossing
representative for each orbit. These two applications use the same local
identity: bounded clock range gives bounded residence time, and specified
clock values give bounded search.

\paragraph{Source and software availability.}
The self-contained supplementary archive includes the retained baseline
implementation, independent checks, local Lean proofs, source hashes,
verification logs, this manuscript source, and build instructions. Run \texttt{python3 reproduce.py} to replay the evidence and
\texttt{make} to typeset the manuscript. The source and release manifests
record SHA256 digests of the supplied files.

\paragraph{AI assistance.}
Generative AI was used in the research conversation to propose and critique
arguments, prepare proof and code drafts, and organize this manuscript.

\begin{thebibliography}{9}
\bibitem{Etse2026}
Kossi Roland Etse.
\newblock How Long Can The Escaping Ant Be Confined?
\newblock In \emph{51st International Symposium on Mathematical Foundations of
Computer Science (MFCS 2026)}, LIPIcs 386, 49:1--49:17, 2026.
\newblock \href{https://doi.org/10.4230/LIPIcs.MFCS.2026.49}{doi:10.4230/LIPIcs.MFCS.2026.49}.

\bibitem{Gajardo2002}
Anah\'{\i} Gajardo, Andr\'{e}s Moreira and Eric Goles.
\newblock Complexity of Langton's ant.
\newblock \emph{Discrete Applied Mathematics} 117:41--50, 2002.
\newblock \href{https://doi.org/10.1016/S0166-218X(00)00334-6}{doi:10.1016/S0166-218X(00)00334-6}.
\end{thebibliography}
\end{document}
